历时数月的毕业设计与论文撰写终于落下帷幕。回首这段充满挑战与收获的开发历程，从最初基于代码逻辑重构执行层架构的迷茫，到最终看着 MCP 审计日志在终端中完美无误地输出，每一次技术突破的背后都离不开诸多师长与亲友的悉心帮助。在此，本人致以最诚挚的谢意。
首先，我要由衷感谢我的指导老师谢小璞、崔笛老师。从论文选题的破局、HIRD 架构的推敲，到系统底层工程实现的纠偏，老师们以其严谨的治学态度和敏锐的技术洞察力，给予了我极其宝贵的指导。
其次，感谢西安电子科技大学软件工程专业四年来的悉心培养。母校扎实严谨的课程体系与浓厚的学术氛围，为我赋予了从算法理论推导到复杂系统编码落地的综合能力。同时，也感谢在课题攻坚期间与我日夜探讨、并肩作战的同窗好友们，是你们的陪伴与技术交流，让原本枯燥的 Debug 过程变得充满乐趣。
最后，我要向我的家人表达最深沉的感恩。感谢你们二十多年如一日的包容、支持与默默付出，是你们给予了我不断向上攀登的精神动力。
未来的学术与工程之路依然漫长，我将带着这份感恩与对技术纯粹的热爱，在生成式 AI 与系统架构的深水区继续笃定前行。